\section{Extraction and additive deletion for lists}
\label{sec:list-transfer}

We first establish a list-coloring form of connectivity extraction. We then
show that deleting a small set together with a suitably sparse set costs only
linearly many colors. Throughout this section, list assignments may have
arbitrary finite palettes.

\begin{lemma}\label{list:elementary}
The following hold for every finite graph $P$.
\begin{enumerate}[label=\textup{(\roman*)}]
\item If $H$ is a subgraph of $P$, then $\ell(H)\le \ell(P)$.
\item For every $Z\subseteq V(P)$,
\[
\ell(P-Z)\ge \ell(P)-|Z|.
\]
\item If $P$ is $q$-degenerate, where $q$ is a nonnegative integer, then
$\ell(P)\le q+1$.
\end{enumerate}
\end{lemma}

\begin{proof}
For the first assertion, extend any list assignment on $H$ with lists of
size at least $\ell(P)$ to the other vertices of $P$, retaining the given
lists on $V(H)$. A list coloring of $P$ restricts to one of $H$; additional
edges among $V(H)$ cause no difficulty.

For the second assertion, write $s=\ell(P-Z)$ and give every vertex of $P$
a list of size at least $s+|Z|$. Greedily color $P[Z]$. At any such step
there are at most $|Z|-1$ colors to avoid. Delete from each remaining list
the colors used on its neighbors in $Z$. At most $|Z|$ entries are removed,
so $P-Z$ can be colored from the remaining lists. This proves
$\ell(P)\le s+|Z|$; the empty-graph cases hold as well under our convention.

For the third assertion, repeatedly remove a vertex of degree at most $q$
from the current graph. In the reverse order, every vertex has at most $q$
already colored neighbors, so lists of size at least $q+1$ suffice.
\end{proof}

We use the constraint-budget method of Gir\~ao and
Narayanan~\cite{GiraoNarayanan2022}, with slack chosen to obtain the
extraction estimate needed here; see also \cite[Section~3.4]{Postle2020}.
The precoloring invariant pays for the
constraints introduced when a separator is colored. A proper list
precoloring on $A$ means a proper coloring of $P[A]$ assigning to each
$v\in A$ a color in its given list.

\begin{lemma}\label{list:constrained-coloring}
Let $k\ge1$ and $m\ge0$ be integers. Suppose every induced $k$-connected
subgraph of $P$ has list chromatic number at most $m$. Give each vertex $v$
a list $\mathcal L(v)$ of size at least $m+6k$. Suppose also that:
\begin{enumerate}[label=\textup{(\roman*)}]
\item a set $A\subseteq V(P)$ has a proper list precoloring $\varphi$;
\item for each $v\notin A$, a finite set $B_v$ of forbidden colors is
specified, with $|B_v|<k$;
\item the constraints satisfy
\begin{equation}\label{list:eq-budget}
k|A|+\sum_{v\notin A}|B_v|\le2k^2.
\end{equation}
\end{enumerate}
Then $P$ has a proper $\mathcal L$-coloring extending $\varphi$ and avoiding
$B_v$ at every $v\notin A$.
\end{lemma}

\begin{proof}
We induct on $|V(P)|$. The hypothesis about induced $k$-connected subgraphs
is inherited by induced subgraphs of $P$. The budget gives $|A|\le2k$.

If $|V(P)|\le2k$, extend $\varphi$ greedily. At each vertex there are at
most $2k-1$ colors used by already colored vertices and at most $k-1$
forbidden colors. Since the lists have size at least $m+6k$, a choice is
available. This includes the empty graph.

If $P$ is $k$-connected, then $\ell(P)\le m$ by hypothesis. For every
$v\notin A$, remove $B_v$ and the colors on its precolored neighbors from
$\mathcal L(v)$. The remaining list has size at least
\[
m+6k-|A|-|B_v|\ge m+3k+1\ge m.
\]
By Lemma~\ref{list:elementary}, $\ell(P-A)\le m$, so these lists color
$P-A$. Combining this coloring with $\varphi$ proves the assertion in
this case. If $P-A$ is empty there is nothing to extend.

It remains to consider $|V(P)|>2k$ when $P$ is not $k$-connected. Choose
a partition
\[
V(P)=V_1\,\dot\cup\,Z\,\dot\cup\,V_2,
\]
where $V_1,V_2$ are nonempty, there are no edges between them, and
$|Z|<k$. Assign to each vertex its constraint cost
\[
w(v)=
\begin{cases}
k,&v\in A,\\
|B_v|,&v\notin A,
\end{cases}
\qquad
w(U)=\sum_{v\in U}w(v).
\]
Interchange the sides if necessary so that $w(V_1)\ge w(V_2)$.
We color the side with the larger constraint cost first, leaving the
other side to absorb the cost of precoloring the whole separator.
In particular,
\begin{equation}\label{list:eq-second-side-cost}
w(V_2)\le k^2.
\end{equation}

We first color $P_1=P[V_1\cup Z]$. Retain the precoloring on
$A_1=A\cap(V_1\cup Z)$. At its other vertices use forbidden sets
\[
B'_v=
\begin{cases}
B_v\cup\{\varphi(x):x\in A\cap V_2,\ xv\in E(P)\},&v\in Z\setminus A,\\
B_v,&v\in V_1\setminus A.
\end{cases}
\]
These extra prohibitions make the eventual separator coloring compatible
with the precolored vertices on the other side. Their total cost obeys
\begin{align}
k|A_1|+\sum_{v\in V(P_1)\setminus A_1}|B'_v|
&\le w(V(P))-w(V_2)+|Z|\,|A\cap V_2|\notag\\
&\le w(V(P))-k|A\cap V_2|+(k-1)|A\cap V_2|\notag\\
&\le2k^2.
\label{list:eq-first-side-cost}
\end{align}
Repeated occurrences of the same color can only reduce the left side.

The sets $B'_v$ need not all have size less than $k$, but each satisfies
\[
|B'_v|\le(k-1)+|A\cap V_2|\le3k-1.
\]
We now convert every vertex with $|B'_v|\ge k$ into a precolored vertex.
Process these vertices one at a time. Choose a color from $\mathcal L(v)$
avoiding $B'_v$ and all colors on its currently precolored neighbors in
$P_1$, and then discard the forbidden set at this newly precolored vertex.
Its contribution to the constraint cost changes from $|B'_v|$ to $k$, so
the total cost does not increase. In particular, throughout this procedure
there are at most $2k$ precolored vertices. Each choice therefore avoids
at most
\[
(3k-1)+2k=5k-1<m+6k
\]
colors and is possible. Every promoted vertex lies in $Z\setminus A$,
since the original forbidden sets on $V_1\setminus A$ had size less than
$k$. Its chosen color also avoids all precolored neighbors in $V_2$, by
the definition of $B'_v$.

After the conversions, the constraints on $P_1$ satisfy all the hypotheses
of the lemma: its precoloring is proper, its remaining forbidden sets have
size less than $k$, and its total cost is at most $2k^2$.
Since $V_2$ is nonempty, $|V(P_1)|<|V(P)|$, and induction supplies a
coloring $\psi_1$ of $P_1$. It respects all the original constraints in
$P_1$. It also avoids every color $\varphi(x)$ at a separator neighbor of
$x\in A\cap V_2$: this was enforced either by an original precoloring,
by a conversion, or by a retained forbidden set.

Next consider $P_2=P[V_2\cup Z]$. Precolor $Z$ according to $\psi_1$ and
retain $\varphi$ on $A\cap V_2$. The resulting precoloring on
$Z\cup(A\cap V_2)$ is proper by the preceding compatibility property.
For each vertex of $V_2\setminus A$, retain its original forbidden set
$B_v$. The new constraint cost is
\[
k|Z|+k|A\cap V_2|+
\sum_{v\in V_2\setminus A}|B_v|
=k|Z|+w(V_2)
\le k(k-1)+k^2<2k^2,
\]
using \eqref{list:eq-second-side-cost}. All its remaining forbidden sets
have size less than $k$. Since $V_1$ is nonempty, induction applies to
$P_2$. Its coloring agrees with $\psi_1$ on $Z$. The two colorings therefore
combine into a proper coloring of $P$, because there are no edges between
$V_1$ and $V_2$. They extend the original precoloring and respect every
original forbidden set.
\end{proof}

\begin{lemma}[List connectivity extraction]\label{lem:list-extraction}
For every integer $k\ge1$, if $\ell(P)>6k$, then some induced
$k$-connected subgraph $H$ of $P$ satisfies
\[
\ell(H)\ge\ell(P)-6k.
\]
\end{lemma}

\begin{proof}
Suppose otherwise and put $m=\ell(P)-6k-1$, a nonnegative integer.
Every induced $k$-connected subgraph then has list chromatic number at
most $m$. Apply Lemma~\ref{list:constrained-coloring} with no precolored
vertices and no forbidden colors. It says that every list assignment of
size at least $m+6k=\ell(P)-1$ is colorable, contradicting the definition
of $\ell(P)$.
\end{proof}

The deletion estimate will reserve colors for a small vertex set. Only
vertices with many neighbors in that set need a probabilistic upper bound
on the number of reserved colors in their lists. Minor exclusion gives a
bound on their number independent of the order of the whole graph.

\begin{lemma}\label{list:high-neighbor-count}
Let $t\ge2$ be an integer, let $h(P)<t$, and let $X\subseteq V(P)$. Then
\[
\bigl|\{v\in V(P)\setminus X:|N_P(v)\cap X|\ge t\}\bigr|
\le\binom{|X|}{2}.
\]
\end{lemma}

\begin{proof}
Process the indicated outside vertices in any order. Assign each one an
as yet unused unordered pair of its neighbors in $X$. Suppose this fails
when processing a vertex $v$. Choose $t$ distinct vertices in
$N_P(v)\cap X$. Every pair of the chosen vertices has already been
assigned to a different previously processed outside vertex adjacent to
both. The corresponding paths of length two have mutually distinct
internal vertices, all outside $X$. They form a subdivision of $K_t$ and
hence give a $K_t$ minor, a contradiction. Thus the assignment never fails
and is an injection into the unordered pairs from $X$.
\end{proof}

\begin{lemma}[Additive list-coloring loss]\label{lem:additive-deletion}
Let $L\ge1$ be the integer constant in Proposition~\ref{prop:small-list}.
For each fixed integer $d\ge0$, put
\[
c_d=4L+d+5.
\]
For every sufficiently large integer $t$, the following holds. Suppose
$h(P)<t$ and $X,Y\subseteq V(P)$ are disjoint sets such that
\begin{enumerate}[label=\textup{(\roman*)}]
\item $|X|\le t^{11/10}$;
\item $P[Y]$ is $dt$-degenerate;
\item every vertex outside $X\cup Y$ has at most $dt$ neighbors in $Y$.
\end{enumerate}
Then
\[
\ell(P-X-Y)\ge\ell(P)-c_dt.
\]
The cutoff is independent of $P$, its order, and all list assignments and
their palettes.
\end{lemma}

\begin{proof}
Write $R=P-X-Y$ and $s=\ell(R)$, and set
\[
u=s+(4L+d+5)t.
\]
We show that $P$ is $u$-choosable. Given any lists of size at least $u$,
first restrict each list to exactly $u$ entries. Denote the resulting
assignment by $\mathcal L$. Its total palette
$\mathcal C=\bigcup_{v\in V(P)}\mathcal L(v)$ is finite, but its size is
otherwise unrestricted.

Independently for each color in $\mathcal C$, reserve that color with
probability
\[
p=\frac{2Lt}{u}.
\]
Since $u\ge(4L+d+5)t>2Lt$, this is a probability strictly between zero and
one. Reservation is global: all occurrences of a color are reserved
together. For each vertex $v$, the number $Z_v$ of reserved entries in
$\mathcal L(v)$ has distribution $\operatorname{Bin}(u,p)$ and mean
$\mu=2Lt$. This mean is independent of $s$, and overlaps between
different lists do not affect this marginal distribution.

For completeness, for every real $\lambda$ its exponential moment obeys
\[
\mathbb E e^{\lambda Z_v}
=(1-p+pe^\lambda)^u
\le\exp\bigl(\mu(e^\lambda-1)\bigr).
\]
Markov's inequality with $\lambda=-\log2$ and $\lambda=\log2$,
respectively, gives
\begin{align}
\Pr(Z_v<\mu/2)
&\le\exp\bigl(-\tfrac{1-\log2}{2}\mu\bigr),
\label{list:eq-lower-tail}\\
\Pr(Z_v>2\mu)
&\le\exp\bigl(-(2\log2-1)\mu\bigr).
\label{list:eq-upper-tail}
\end{align}

Let
\[
T=\{v\in V(P)\setminus X:|N_P(v)\cap X|\ge t\}.
\]
By Lemma~\ref{list:high-neighbor-count}, $|T|\le\binom{|X|}{2}$. Put
$\eta=\min\{(1-\log2)/2,\,2\log2-1\}>0$. A union bound using
\eqref{list:eq-lower-tail} and \eqref{list:eq-upper-tail} shows that the
probability of either of the failures
\[
Z_x<Lt\quad\text{for some }x\in X,
\qquad
Z_v>4Lt\quad\text{for some }v\in T
\]
is at most
\[
(|X|+|T|)e^{-2\eta Lt}
\le2t^{11/5}e^{-2\eta Lt}<1
\]
for sufficiently large $t$. There are only polynomially many tested
vertices even if $P$ has arbitrarily many vertices. No independence
between different vertices' failure events is required. Fix a reservation
for which neither failure occurs.

Use the reserved entries to color $P[X]$. They provide at least $Lt$
colors at each of its vertices, and Proposition~\ref{prop:small-list}
applies because $|X|\le t^{11/10}$ and $h(P[X])<t$. If $X$ is empty this
step is vacuous. Write the coloring as $\varphi_X$.

For each vertex $v\notin X$, delete from its list the colors on its
neighbors in $X$:
\[
\mathcal L_1(v)=\mathcal L(v)\setminus
\{\varphi_X(x):x\in N_P(v)\cap X\}.
\]
If $v\in T$, every deleted entry is reserved, so at most $4Lt$ entries
are deleted. If $v\notin T$, fewer than $t$ entries are deleted because
$v$ has fewer than $t$ neighbors in $X$. In particular, the uniform bound
\begin{equation}\label{list:eq-first-pruning}
|\mathcal L_1(v)|\ge u-(4L+1)t=s+(d+4)t
\end{equation}
holds. We have removed only colors actually used on neighbors in $X$;
there is no need to discard all reserved entries at the other vertices.

The bound in \eqref{list:eq-first-pruning} is at least $dt+1$. By
Lemma~\ref{list:elementary}, the $dt$-degenerate graph $P[Y]$ therefore
has a coloring $\varphi_Y$ from these lists. This coloring has no conflict
with $\varphi_X$.

Finally, at each vertex of $R$ delete from $\mathcal L_1(v)$ the colors
used by its neighbors in $Y$. There are at most $dt$ such neighbors, so
its resulting list has size at least
\[
u-(4L+1)t-dt=s+4t\ge s.
\]
The definition of $s=\ell(R)$ supplies a coloring of $R$ from these
lists, with nothing to do if $R$ is empty. The three colorings combine
properly, since both pruning steps excluded colors used on neighbors in
the previously colored parts. We have colored every $u$-list assignment
on $P$, proving
\[
\ell(P)\le u=\ell(P-X-Y)+c_dt.
\]
Only the fixed cutoff in Proposition~\ref{prop:small-list} and the last
exponential-versus-polynomial estimate impose a sufficiently-large-$t$
requirement. Neither depends on the graph order or palette size.
\end{proof}
